% GENERATED FILE. Edit canonical lessons, then run npm run generate:books.
% canonical-source-hash: fnv1a64:879b3e7486ffb8b8
% canonical-lessons: TA-C449-elam, TA-C449-kuttakai, TA-C449-parivarttanai, TA-C449-mutalitu, TA-C449-ataku-vai

\chapter{Auctions and Leases}
\label{ch:ta-auctions-and-leases}
\hlchaptermodality{449}

\begin{chapteropening}
\textbf{By the end of this chapter:} \emph{I can talk about auctions, leases, transactions and investments, and pawning jewellery.}

Complete the last lesson of chapter 449: I can talk about auctions, leases, transactions and investments, and pawning jewellery.
\end{chapteropening}

\section[\texorpdfstring{\ta{ஏலம்} (ēlam)}{ஏலம் (ēlam)}]{\ta{ஏலம்} (ēlam) — an auction}

\label{lesson:TA-C449-elam}



\begin{quote}
Before the new one: say the Tamil for to be left over, then the Tamil for income.
\end{quote}

\subsection*{You'll want to know: \ta{ஏலம்}}
\textbf{\ta{ஏலம்}} — \emph{ēlam} — \textquotedblleft{}an auction\textquotedblright{}.

\begin{quote}
\ta{பழைய} \ta{கார்களை} \ta{ஏலத்தில்} \ta{விற்றார்கள்}.
\end{quote}

\begin{quote}
\emph{paḻaiya kārkaḷai ēlattil viṟṟārkaḷ.}
\end{quote}

\begin{quote}
They sold the old cars at an auction.
\end{quote}

\ta{ஏலம்} \ta{விடு} is to put up for auction.

\begin{cousinweb}
\textbf{\ta{ஏலம்}} (\emph{ēlam}) is usually traced to Portuguese \emph{leilão}, \textquotedblleft{}an auction\textquotedblright{}.
\end{cousinweb}

\begin{grammarlens}[title={What you've built}]
One more word for getting around.
\end{grammarlens}

\subsection*{Your turn}
\begin{itemize}
\raggedright
  \item \emph{Say it:} \emph{ēlam}
  \item \emph{Say it:} \emph{ēlam}, once more
  \item \emph{Say it:} \emph{paḻaiya kārkaḷai ēlattil viṟṟārkaḷ.} — They sold the old cars at an auction.
  \item \emph{Recall:} say the Tamil for to be left over, then the Tamil for income, then say \emph{ēlam} again
\end{itemize}

\begin{tcolorbox}[breakable,colback=teal!4,colframe=teal!35!black,title={Before you move on}]
What does \ta{ஏலம்} mean? (\textquotedblleft{}An auction\textquotedblright{}.) Say it once more.
\end{tcolorbox}
\section[\texorpdfstring{\ta{குத்தகை} (kuttakai)}{குத்தகை (kuttakai)}]{\ta{குத்தகை} (kuttakai) — a lease; a tenancy}

\label{lesson:TA-C449-kuttakai}



\begin{quote}
Before the new one: say the Tamil for income, then the Tamil for an auction.
\end{quote}

\subsection*{You'll want to know: \ta{குத்தகை}}
\textbf{\ta{குத்தகை}} — \emph{kuttakai} — \textquotedblleft{}a lease; a tenancy\textquotedblright{}.

\begin{quote}
\ta{அவர்} \ta{அந்த} \ta{நிலத்தைக்} \ta{குத்தகைக்கு} \ta{எடுத்தார்}.
\end{quote}

\begin{quote}
\emph{avar anta nilattaik kuttakaikku eṭuttār.}
\end{quote}

\begin{quote}
He took that land on lease.
\end{quote}

\ta{குத்தகைக்கு} \ta{எடு} is to take on lease.

\begin{cousinweb}
\textbf{\ta{குத்தகை}} (\emph{kuttakai}), \textquotedblleft{}a lease, a contract\textquotedblright{}, is a word whose origin is not settled.
\end{cousinweb}

\begin{grammarlens}[title={What you've built}]
One more word for getting around.
\end{grammarlens}

\subsection*{Your turn}
\begin{itemize}
\raggedright
  \item \emph{Say it:} \emph{kuttakai}
  \item \emph{Say it:} \emph{kuttakai}, once more
  \item \emph{Say it:} \emph{avar anta nilattaik kuttakaikku eṭuttār.} — He took that land on lease.
  \item \emph{Recall:} say the Tamil for income, then the Tamil for an auction, then say \emph{kuttakai} again
\end{itemize}

\begin{tcolorbox}[breakable,colback=teal!4,colframe=teal!35!black,title={Before you move on}]
What does \ta{குத்தகை} mean? (\textquotedblleft{}A lease; a tenancy\textquotedblright{}.) Say it once more.
\end{tcolorbox}
\section[\texorpdfstring{\ta{பரிவர்த்தனை} (parivarttaṉai)}{பரிவர்த்தனை (parivarttaṉai)}]{\ta{பரிவர்த்தனை} (parivarttaṉai) — a transaction; an exchange}

\label{lesson:TA-C449-parivarttanai}



\begin{quote}
Before the new one: say the Tamil for an auction, then the Tamil for a lease.
\end{quote}

\subsection*{You'll want to know: \ta{பரிவர்த்தனை}}
\textbf{\ta{பரிவர்த்தனை}} — \emph{parivarttaṉai} — \textquotedblleft{}a transaction; an exchange\textquotedblright{}.

\begin{quote}
\ta{இணையத்தில்} \ta{பணப்} \ta{பரிவர்த்தனை} \ta{செய்வது} \ta{எளிது}.
\end{quote}

\begin{quote}
\emph{iṇaiyattil paṇap parivarttaṉai ceyvatu eḷitu.}
\end{quote}

\begin{quote}
Making money transactions on the internet is easy.
\end{quote}

\ta{பணப்} \ta{பரிவர்த்தனை} is a money transaction.

\begin{cousinweb}
\textbf{\ta{பரிவர்த்தனை}} (\emph{parivarttaṉai}) is borrowed from Sanskrit \emph{parivartana}, \textquotedblleft{}turning round, moving to and fro\textquotedblright{} (Monier-Williams), used for exchange.
\end{cousinweb}

\begin{grammarlens}[title={What you've built}]
One more word for getting around.
\end{grammarlens}

\subsection*{Your turn}
\begin{itemize}
\raggedright
  \item \emph{Say it:} \emph{parivarttaṉai}
  \item \emph{Say it:} \emph{parivarttaṉai}, once more
  \item \emph{Say it:} \emph{iṇaiyattil paṇap parivarttaṉai ceyvatu eḷitu.} — Making money transactions on the internet is easy.
  \item \emph{Recall:} say the Tamil for an auction, then the Tamil for a lease, then say \emph{parivarttaṉai} again
\end{itemize}

\begin{tcolorbox}[breakable,colback=teal!4,colframe=teal!35!black,title={Before you move on}]
What does \ta{பரிவர்த்தனை} mean? (\textquotedblleft{}A transaction; an exchange\textquotedblright{}.) Say it once more.
\end{tcolorbox}
\section[\texorpdfstring{\ta{முதலீடு} (mutalīṭu)}{முதலீடு (mutalīṭu)}]{\ta{முதலீடு} (mutalīṭu) — an investment}

\label{lesson:TA-C449-mutalitu}



\begin{quote}
Before the new one: say the Tamil for a lease, then the Tamil for a transaction.
\end{quote}

\subsection*{You'll want to know: \ta{முதலீடு}}
\textbf{\ta{முதலீடு}} — \emph{mutalīṭu} — \textquotedblleft{}an investment\textquotedblright{}.

\begin{quote}
\ta{அவர்} \ta{தன்} \ta{பணத்தை} \ta{நிலத்தில்} \ta{முதலீடு} \ta{செய்தார்}.
\end{quote}

\begin{quote}
\emph{avar taṉ paṇattai nilattil mutalīṭu ceytār.}
\end{quote}

\begin{quote}
He invested his money in land.
\end{quote}

\ta{முதலீடு} \ta{செய்} is to invest.

\begin{cousinweb}
\textbf{\ta{முதலீடு}} (\emph{mutalīṭu}) is built from \ta{முதல்}, \textquotedblleft{}the first, the principal, capital\textquotedblright{} (DEDR 4950), and \emph{īṭu}, \textquotedblleft{}putting in\textquotedblright{} (DEDR 442).
\end{cousinweb}

\begin{grammarlens}[title={What you've built}]
One more word for getting around.
\end{grammarlens}

\subsection*{Your turn}
\begin{itemize}
\raggedright
  \item \emph{Say it:} \emph{mutalīṭu}
  \item \emph{Say it:} \emph{mutalīṭu}, once more
  \item \emph{Say it:} \emph{avar taṉ paṇattai nilattil mutalīṭu ceytār.} — He invested his money in land.
  \item \emph{Recall:} say the Tamil for a lease, then the Tamil for a transaction, then say \emph{mutalīṭu} again
\end{itemize}

\begin{tcolorbox}[breakable,colback=teal!4,colframe=teal!35!black,title={Before you move on}]
What does \ta{முதலீடு} mean? (\textquotedblleft{}An investment\textquotedblright{}.) Say it once more.
\end{tcolorbox}
\section[\texorpdfstring{\ta{அடகு} \ta{வை} (aṭaku vai)}{அடகு வை (aṭaku vai)}]{\ta{அடகு} \ta{வை} (aṭaku vai) — to pawn; to pledge}

\label{lesson:TA-C449-ataku-vai}



\begin{quote}
Before the new one: say the Tamil for a transaction, then the Tamil for an investment.
\end{quote}

\subsection*{You'll want to know: \ta{அடகு} \ta{வை}}
\textbf{\ta{அடகு} \ta{வை}} — \emph{aṭaku vai} — \textquotedblleft{}to pawn; to pledge\textquotedblright{}.

\begin{quote}
\ta{அவள்} \ta{தன்} \ta{வளையல்களை} \ta{அடகு} \ta{வைத்தாள்}.
\end{quote}

\begin{quote}
\emph{avaḷ taṉ vaḷaiyalkaḷai aṭaku vaittāḷ.}
\end{quote}

\begin{quote}
She pawned her bangles.
\end{quote}

Past (I …): \ta{அடகு} \ta{வைத்தேன்}, \emph{aṭaku vaittēṉ}.

\ta{அடகுக்} \ta{கடை} is a pawnshop.

\begin{cousinweb}
\textbf{\ta{அடகு} \ta{வை}} (\emph{aṭaku vai}) joins \ta{அடகு}, \textquotedblleft{}a pledge, mortgaged property\textquotedblright{}, which is native Dravidian (DEDR 79), to \ta{வை}, \textquotedblleft{}to put\textquotedblright{}.
\end{cousinweb}

\begin{grammarlens}[title={What you've built}]
One more word for getting around.
\end{grammarlens}

\subsection*{Your turn}
\begin{itemize}
\raggedright
  \item \emph{Say it:} \emph{aṭaku vai}
  \item \emph{Say it:} \emph{aṭaku vai}, once more
  \item \emph{Say it:} \emph{avaḷ taṉ vaḷaiyalkaḷai aṭaku vaittāḷ.} — She pawned her bangles.
  \item \emph{Recall:} say the Tamil for a transaction, then the Tamil for an investment, then say \emph{aṭaku vai} again
\end{itemize}

\begin{tcolorbox}[breakable,colback=teal!4,colframe=teal!35!black,title={Before you move on}]
What does \ta{அடகு} \ta{வை} mean? (\textquotedblleft{}To pawn; to pledge\textquotedblright{}.) Say it once more.
\end{tcolorbox}
